\section{Related Work}
\label{sec:related}

\paragraph{Shielding and recoverability-set safe RL.}
Shielding~\cite{alshiekh2018safe} synthesizes a finite-state shield from a temporal-logic specification that intercepts policy outputs at runtime; Model-Predictive Shielding~\cite{bastani2019mps}, Recovery RL~\cite{thananjeyan2021recovery}, and RL power-grid shielding~\cite{hrlshield2026} relax binary terminal gates by admitting recoverability-set states or screening actions with forward simulation. All assume a fixed action space and require either a backup/recovery policy or a dynamics-model lookahead. \silr{} instead shields \emph{black-box LLM tool-call dispatch} with neither: proposals are admitted by single-step shadow simulation under a product order on the branch-level violation state, at constant per-step cost. Where these lineages apply a fixed admissibility filter to a known action model, \silr{} makes admission itself the design variable---a graded predicate over violation geometry.

\paragraph{Runtime guards for LLM agents.}
Recent work adds runtime enforcement to tool-using LLM agents through architectural gates, guard agents, or rule languages: SARC's Pre-Action Gate~\cite{sarc2026} (the closest binary baseline, admitting a tool call only when a hand-written predicate holds), AgentSpec's runtime-enforcement DSL~\cite{agentspec2026}, GuardAgent~\cite{guardagent2025} and ShieldAgent~\cite{shieldagent2025} (compiling guard policies), and VeriGuard~\cite{veriguard2025} (verified code generation plus monitoring). These make LLM agents more governable but do not target post-violation recovery admission: a terminal predicate (as a SARC-style Pre-Action Gate instantiates) preserves safety but deadlocks multi-step recovery (Sec.~\ref{sec:eval:rq1}).

\paragraph{LLM agents for power-grid control.}
Grid-Agent~\cite{gridagent2025} is the closest same-domain near-neighbor: it couples LLM agents with sandboxed power-flow validation, rollback of negative actions, and monotonic-progress behavior for power-grid control. The mechanisms differ in kind: Grid-Agent's check is a \emph{post-hoc} rollback that reverts an action unless the violation set is resolved without introducing new violations---a support-level criterion with no per-branch severity guard---whereas \silr{}'s product order is a \emph{pre-action} admission predicate that never applies a geometry-poor step and emits graded verdicts reusable as a training signal.

\paragraph{Multi-objective optimization and scalarization.}
Scalarization trade-offs are classically studied as a \emph{static} representational question---e.g., linear scalarization cannot reach non-convex Pareto regions~\cite{miettinen1999nonlinear}. The scalar projection trap (Sec.~\ref{sec:method:trap}) is a \emph{behavioral} phenomenon of the runtime gate: in a ReAct loop the gate is a commitment operator, so an admitted aggregate-improving step restarts the per-step search from a plateau state---a failure of the gate's trajectory dynamics, not merely of an aggregate's expressiveness.

\paragraph{Adversarial LLM defenses.}
Prompt injection~\cite{greshake2023injection} is a documented attack vector; prior defenses sanitize inputs or outputs. \silr{} adopts a third stance: the LLM \emph{may} be compromised, and safety is enforced after the proposal by shadow-executing it on the simulator, shifting the trust boundary from the LLM to the verifier and simulator.

\paragraph{Positioning.}
Table~\ref{tab:related} places \silr{} against the closest systems along four axes: \emph{graded} admission (beyond binary accept/reject), \emph{retained multi-component violation geometry} (vs.\ a scalar aggregate), \emph{black-box LLM tool proposals} (no fixed action space), and \emph{neither a backup/recovery policy nor a multi-step dynamics lookahead}. \silr{} is the only point meeting all four---this combination is the contribution.

\begin{table}[tb]
\centering
\caption{Positioning of \silr{} against the closest runtime-safety systems. ``Geometry'': retains the (support, severity) violation state vs.\ a scalar aggregate; ``No lookahead/backup'': single-step shadow only, no recovery policy or dynamics rollout.}
\label{tab:related}
\footnotesize
\setlength{\tabcolsep}{4pt}
\begin{tabular*}{\textwidth}{@{\extracolsep{\fill}}lcccc}
\toprule
System & Graded & Geometry & \makecell{Black-box\\LLM} & \makecell{No lookahead/\\backup} \\
\midrule
SARC Pre-Action Gate~\cite{sarc2026}        & \ding{55} & \ding{55} & \ding{51} & \ding{51} \\
AgentSpec DSL~\cite{agentspec2026}          & \ding{55} & \ding{55} & \ding{51} & \ding{51} \\
GuardAgent/ShieldAgent~\cite{guardagent2025,shieldagent2025} & \ding{55} & \ding{55} & \ding{51} & \ding{51} \\
Grid-Agent~\cite{gridagent2025}             & \ding{55} & \ding{55}$^\dagger$ & \ding{51} & \ding{55}$^\ddagger$ \\
MPS / Recovery RL~\cite{bastani2019mps,thananjeyan2021recovery} & \ding{55} & \ding{55} & \ding{55} & \ding{55} \\
\silr{} (ours)                              & \ding{51} & \ding{51} & \ding{51} & \ding{51} \\
\bottomrule
\end{tabular*}\\[2pt]
{\footnotesize $^\dagger$support-level violation validation, no per-branch severity guard; $^\ddagger$applies then reverts (post-hoc rollback).}
\end{table}
